\section{Експериментални резултати и анализ}
\label{sec:results}

Разделът представя резултатите в следния ред: степен на съвпадение, представителни съпоставки,
систематизация на нарушенията на независимостта, цена на реализацията и ограничения. Основното
съдържание е в §\ref{subsec:taxonomy}: броят съвпадащи интерфейси е състояние на конкретна
реализация към конкретна дата, докато систематизацията на \emph{къде} и \emph{защо} абстракцията
протича се обобщава извън артефакта.

\subsection{Степен на съвпадение}
\label{subsec:agreement}

Обобщението е в \tabref{tab:parity}. От \boardPages{} интерфейса на всяка от четирите платформи ---
iOS, Mac Catalyst, Android и Windows --- броят на съществено различаващите се по входния път
\emph{по код} е \boardRedios{}, \boardRedmaccatalyst{}, \boardRedandroid{} и \boardRedwindows{}:
бордът е практически изчистен от съществени разлики на четирите платформи едновременно. Единственото
изключение е по входния път \emph{от описанието} на iOS --- една двойка, разгледана в
§\ref{subsec:visual}. Пълната картина е топлинна карта в Приложение~В \tabref{tab:heat}, всяка
клетка носи измереното структурно сходство, а цветът следва стойността (зелено 1{,}000 --- жълто
0{,}75 --- червено 0{,}50) --- изключенията се намират с поглед.

\input{generated/board_tables}

\subsubsection{Интерфейси с побайтово съвпадение}
\label{subsec:integrity_finding}

Част от интерфейсите се изобразяват от системата \emph{побайтово идентично} с еталона --- не
близко, а един и същ файл по съдържание \tabref{tab:byteid}.

\input{generated/byteid_table}

Резултатът заслужава проверка, преди да бъде приет: побайтово съвпадение е и признакът на
известен режим на отказ --- кадър, попаднал под две имена (§\ref{subsec:integrity}).
Разграничението е направено с три независими проверки:

\begin{enumerate}[topsep=2pt,itemsep=1pt,leftmargin=*]
\item \textbf{Времена на файловете.} Трите колони носят различни времена, подредени по реда на
      заснемащите проходи --- дни и часове една от друга; размножен кадър би носил един момент.
\item \textbf{Темите остават различни.} Във всичките 172 случая светлото и тъмното заснемане се
      различават; известният режим на отказ има точно обратния признак.
\item \textbf{Съгласуваност с останалите платформи.} Платформата, при която заснемането обхваща
      прозорец с жива системна лента и сянка, има \emph{нула} такива случая --- очакваното, защото
      там побайтова устойчивост е невъзможна по начина на изграждане.
\end{enumerate}

Постижимо е тук, защото двете колони минават през един и същ заснемащ инструмент, при една и съща
разделителна способност, с фиксиран часовник в лентата и детерминистичен кодер: еднакви пиксели
дават еднакви байтове.

Числата в \tabref{tab:byteid} са \emph{по-строг}, не по-слаб критерий спрямо структурно сходство
1{,}0000: пикселът се брои различаващ се едва след разлика над 25 от 255 по канал
(§\ref{subsec:scoring}), за да не наказва заглаждането на шрифта, а самото структурно сходство се
показва с четири знака (,,1{,}0000“ означава $\geq 0{,}99995$, не точно единица) --- затова кадър
с оценка ,,точно съвпадение“ невинаги е и байт по байт същият файл; двете мерки отговарят на
различни въпроси. Системната лента участва в оценяването навсякъде освен на едната мобилна
платформа, където горните 140 реда се отрязват (§\ref{subsec:integrity}); на настолната платформа
тя носи \emph{заглавието на прозореца} --- различно по замисъл между еталон и система --- и остава
източник на дребни, немаскируеми разлики, независимо от съдържанието отдолу.

\textbf{Следствие за инструмента, не за данните.} Механичната проверка от §\ref{subsec:integrity}
третира всяко междуколонно съвпадение като дефект --- вярно, докато системата не достигаше точно
съвпадение, но източник на лъжливи тревоги, чийто брой расте точно с качеството на реализацията.
Безусловно вярно остава по-тясното условие: междуколонно съвпадение \emph{в съчетание със}
съвпадение между темите. Критерий за достоверност може да остарее заради успех на измерваната
система --- основание той да се преразглежда, не да се приема за вечен.

\subsubsection{Платформи с ограничена съпоставимост}

Две платформи носят ограничения, отчитани отделно, не осреднени с останалите:

\begin{enumerate}[topsep=2pt,itemsep=1pt,leftmargin=*]
\item \textbf{Тъмната тема на една от мобилните платформи не е съпоставима.} Еталонът там
      изобразява светъл интерфейс независимо от поисканата тема (средна яркост на тялото 137,7
      срещу 139,3 --- на практика еднакви), докато системата изобразява действително тъмен;
      системата е права срещу неработещ еталон, но точно затова числото не носи информация.
\item \textbf{Платформата с различния native инструментариум} не се оценява количествено по замисъл
      (§\ref{subsec:exp_setup}) --- за нея се отчитат пълнота на съдържанието и съгласуваност между
      двата входни пътя.
\end{enumerate}

\subsubsection{Динамика на резултата}

Дневникът на разработката записва състоянието на борда след всяка кампания за отстраняване на
разминавания: резултатът е достигнат постепенно, а не наведнъж --- само за последните девет дни от
работата бяха направени 111 отделни ангажимента по заснемащата и сравняващата инфраструктура, всеки
затварящ по едно установено разминаване. Записите са свободен текст, не машинно четими данни, затова
тук не се привежда графика --- възстановена крива от непълни записи би внушила точност, каквато
източникът няма; систематизацията в §\ref{subsec:taxonomy} е носителят на този опит.

\subsection{Представителни съпоставки}
\label{subsec:visual}

Фигурите по-долу поставят еталона и системата една до друга при идентично входно описание.
Изображенията се вземат чрез символни връзки към борда (\code{tools/gen\_tex.py}), тоест са
\emph{същите файлове}, които са били оценени --- не подбрани отделно за изложението.

\evidencepair{maccatalyst}{absolute_layout}{light}{Абсолютно разположение, светла тема ---
контрола, при която разминаването би било изместване на всеки елемент поотделно.}

\evidencepair{maccatalyst}{box_view}{dark}{Правоъгълници с цвят, градиент и закръгление, тъмна
тема --- проверява цвят, геометрия и отсичане едновременно.}

\evidencepair{maccatalyst}{label}{light}{Текстова контрола --- проверява шрифт, разстояние между
знаците и водещата линия едновременно.}

\evidencepair{maccatalyst}{basic_grouping}{light}{Списък с групиране --- най-натовареният клас
интерфейс: заглавия на групи, вътрешни отстъпи и превъртане.}

Двойка ПРЕДИ и СЛЕД отстраняване на конкретно разминаване не се привежда като фигура: предишните
състояния не са запазени като заснемания, а само като измерени стойности в дневника (например
,,6 px / 10,5\,\% $\rightarrow$ 0 / 0,5\,\%“). Числата са дадени при съответния клас в
§\ref{subsec:taxonomy}; изображенията не могат да бъдат възстановени, без съответната версия да
бъде компилирана наново.

За анимираните интерфейси (\gifPageCount{} на брой) фигурата би показала \emph{първия кадър};
където такъв се използва, това се отбелязва в легендата.

\subsubsection{Най-слабо съвпадащите интерфейси}

Прозрачността изисква да бъдат показани не само добрите случаи. Таблицата се генерира от данните;
само колоната ,,причина“ се поддържа на ръка, защото оценката е измерване, а причината е находка
--- интерфейс без установена причина се отпечатва като ,,открит“, а не получава правдоподобно
обяснение.

\input{generated/worst_table}

Към момента на писане таблицата съдържа един-единствен ред: същата двойка (iOS, входен път от
описанието), която носи и единствената ,,съществена“ клетка в целия борд (§\ref{subsec:agreement}).
Разследването ѝ (§\ref{subsec:motion_scoring}) я класифицира като шум при кацане на анимация,
установен чрез контрола ,,еталонът срещу себе си“: повторни стартирания на самия еталон се
разминават помежду си повече, отколкото системата с еталона. Колоната ,,причина“ отбелязва
,,открит“, не измислено обяснение: дали прагът за допустимо изместване при кацане да бъде вдигнат
--- или платформата изключена от изискването за възпроизводимо кацане --- е оставено съзнателно
отворено.

По-ранни разследвания оставиха и трайни методологични изводи. Интерфейсът \code{image} показа защо
двата входни пътя не се събират в едно число: на Android само пътят по код се разминаваше съществено
(0{,}478), докато пътят от описанието изобразяваше вярно (0{,}999) --- осреднени, случаят би
изчезнал безследно. Точно затова §\ref{subsec:rulings} изисква четирите сравнения да се отчитат
поотделно, не като едно число на интерфейс.

\evidencepair{ios}{image}{light}{\code{image} на iOS --- илюстрация на находката за двата входни
пътя: в колоната на системата липсва само изтегляният по мрежата образ, останалите източници се
изобразяват вярно.}

\evidencepair{windows}{context_flyout}{light}{Интерфейс, изключен от количествената оценка по
правилото за артефакт на измерването (§\ref{subsec:rulings}): вгражда жива уеб страница, изтеглена
от двете колони в различни моменти.}

\subsection{Систематизация на нарушенията на независимостта}
\label{subsec:taxonomy}

Основният резултат на работата. Всяко разминаване, установено в хода на измерването и доведено до
първопричина, попада в един от изброените по-долу класове. Систематизацията не е съставена по
интуиция: тя обобщава \emph{доведени до първопричина} случаи от кампаниите за отстраняване на
разминавания, всеки документиран с измерване преди и след.

Общата констатация, която следва от разпределението им, е нетривиална и противоречи на очакването:
\emph{абстракцията не протича на границата на програмния интерфейс на контролите}. Съответствието
,,свойство\,$\rightarrow$\,native свойство“ се пренася почти без съпротива. Протичането е на
местата, където всеки инструментариум кодира \term{неизказана уговорка} --- допускане, което не е
част от документирания му интерфейс, но е част от поведението му.

\begin{table}[H]
\caption{Класове нарушения на независимостта}
\label{tab:taxonomy}
\centering\small
\begin{tabular}{@{}L{0.7cm}L{5.2cm}L{4.0cm}L{3.6cm}@{}}
\hline
\textbf{№} & \textbf{Клас} & \textbf{Къде се проявява} & \textbf{Откриваем без екран?} \\
\hline
1 & допускания при наслагването & състав от слоеве & не \\
2 & геометрични конвенции & измерване/подреждане & \emph{да} \\
3 & разрешаване на стойности и ред & таблици, стилове & \emph{да} \\
4 & входове от средата & всички платформи наведнъж & не \\
5 & семантика на отстъпите & ръбове, вложени изгледи & частично \\
6 & поведение само от реализацията & разклонения по тип & \emph{да} \\
7 & не са дефекти & --- & --- \\
\hline
\end{tabular}
\end{table}

Броят доведени до първопричина случаи във всеки клас не се привежда като число: той би отразявал
реда, в който дефектите са били търсени, а не действителното им разпределение. Информативно е
съотношението в последната колона --- три от шестте действителни класа са откриваеми без екран,
три изискват сравнение на изображения.

\subsubsection{Клас 1 --- допускания при наслагването}

Разминавания, при които системата рисува вярно, но резултатът се наслагва върху основа с други
свойства от очакваните --- изолираното измерване на слоя показва правилен резултат, грешката се
проявява едва в състава.

\textbf{Разгледан пример --- фон на прозореца, изграден от два слоя.} Еталонът установява на
прозореца системен ефект зад съдържанието (полупрозрачно замъгляване на работния плот), а върху
него --- \emph{полупрозрачен} слой за областта на съдържанието. Наблюдаваният цвят на тялото е
\emph{съставен} от двете. Системата възпроизведе първо само първия слой: измереното --- стойност
232 в светла тема --- вече доказваше, че ефектът работи (непрозрачно запълване по подразбиране би
дало 243), но цялото не съвпадаше, защото липсваше вторият, полупрозрачен слой; добавянето на
\emph{непрозрачно} запълване на основата, обратно, се наслагва \emph{пред} системния ефект и го
обезсилва напълно. Правилното решение е основата да не се боядисва изобщо там, където ефектът е
наличен.

Класът е поучителен по два начина: изолираното измерване на всеки слой поотделно дава правилен
резултат --- грешката е само в състава; и посоката на поправката е обратна на интуицията --- не
,,добави още боя“, а ,,махни я, за да се вижда това, което е отдолу“. (Тъмната тема отчиташе яркост
32 при същото положение; заснемане отпреди поправката не е запазено --- бордът съхранява текущото
състояние, не историята на изображенията --- запазени са само измерените стойности.)

\subsubsection{Клас 2 --- геометрични конвенции}

Класът с най-много случаи: всяка от двете страни има вътрешно съгласувана конвенция, конвенциите не
съвпадат, и нито една не е записана в интерфейса.

\textbf{Разгледан пример --- несиметрично отчитане на външния отстъп.} Договорът за измерване
изисква измерващата фаза да \emph{прибави} външния отстъп към съобщения размер, а подреждащата ---
да го \emph{извади} от рамката; при нулев отстъп двете се компенсират, затова грешката е невидима
в най-простия случай. В реализацията базовият изглед изпълняваше и двете, но преопределянето за
контейнерите --- само изваждането. Следствието е, че рамката на \emph{всеки} контейнер губеше
двукратния си отстъп: при поискана ширина 100 и отстъп 12 подреждането даваше 76 вместо 100.

Двете фази и посоката на прибавяне/изваждане са показани на \figref{fig:measure_arrange}.

Преди установяването на действителната причина бяха проверени и \emph{отхвърлени с измерване} две
правдоподобни хипотези --- че вътрешният отстъп свива изричния размер, и че конкретна контрола
пренебрегва собствения си външен отстъп. Отхвърлени чрез насочена проверка без екран, не чрез
разсъждение, те са записани, за да не бъдат изпробвани отново.

\textbf{Втори пример --- половинпикселно свиване на контура.} Инструментариумът за чертане свива
контура на фигура с половин единица навътре, за да го центрира върху ръба. Същата фигура обаче служи
и като \emph{път на отсичане} за съдържанието вътре в нея, където свиването няма основание --- води
до отрязване на половин пиксел по целия периметър, дребно на всеки ръб и видимо като цвят по цялата
рамка.

\textbf{Трети пример --- координатната система при подреждане} (§\ref{subsec:layout_design}):
абсолютни спрямо относителни на контейнера координати. Погрешният избор води до двойно отместване
само при вложен контейнер с ненулево начало --- не при най-простите интерфейси, типично за класа.

\subsubsection{Клас 3 --- разрешаване на стойности и ред на прилагане}

\textbf{Разгледан пример --- зависимост от реда в таблицата за преобразуване.} Размерът на бутон се
изчислява от вътрешните отстъпи на native контролата, които трябва да включват дебелината на
контура. Двете свойства се прилагат от таблицата в определен ред --- отстъп преди контур --- така
че при изчисляване на размера дебелината още е нула и контурът никога не достига до отстъпите.
Еталонът оцелява при същия ред, защото на по-горното си ниво \emph{замества} действието за отстъп с
преизчисляване при измерване, тоест чете дебелината в момента на измерването: \emph{редът на
прилагане в таблицата е част от наблюдаваното поведение}, не вътрешна подробност. Решението е
изчислението да получава дебелината като аргумент, независимо от реда --- поправка на класа, не на
случая.

\textbf{Втори пример --- минимална ширина, наложена от стила на платформата.} С две присвоявания
еталонът занулява минималната ширина и височина на превключвател без текст. На една от
платформите те са \emph{безрезултатни}: действието за минимална ширина в таблицата изтрива локално
зададената стойност, след което влиза минимумът от стила на инструментариума --- 120 единици,
надделяващ дори над изрично зададена от приложението ширина. Системата пренесе двата реда буквално и
получи 32 единици вместо 120. Вярното пренасяне е \emph{да не се пренасят}: незадаването е точно
състоянието след изтриването, и минимумът от стила се проявява от само себе си --- буквалното
пренасяне на ред от еталона тук \emph{разминава} поведението.

\textbf{Трети пример --- ,,незададено“ срещу ,,зададено на стойността по подразбиране“.} Проверка,
която пита за \emph{стойността} вместо за \emph{наличието} на стойност, не разграничава двата
случая: контрола с изрично зададен черен цвят на текста се третира като без зададен цвят и получава
цвета по подразбиране на платформата --- бял при тъмна тема, текстът невидим. Поправката е изрична
проверка ,,това свойство задавано ли е“, различна от сравнение със стойността.

\subsubsection{Клас 4 --- входове от средата}

Разминавания, чийто източник е стойност, която приложението \emph{не} задава, а получава от
средата: проявяват се едновременно на всички платформи, а причината е на едно място.

\textbf{Разгледан пример --- източникът на темата.} Приложението има слот за тема, но началната му
стойност \emph{не} е ,,светла“ --- еталонът я установява в конструктора си от операционната система.
Системата отначало пропускаше този един ред, при което слотът оставаше на подразбиращата се
стойност: на машина с тъмна тема всяка платформа едновременно показва интерфейс, изграден по светли
правила, върху тъмна системна основа. Клас, при който всяка платформа се проявява поотделно, почти
никога не е от този вид.

Интерфейсът, който свързва темата (\code{app\_theme\_binding}), е и този, на който разминаването се
вижда най-пряко. Измереното след поправката: от 4,6\,\% различаващи се пиксела до 0,08\,\% --- тоест
от съществено разминаване до точно съвпадение. Заснемане отпреди поправката не е запазено по същата
причина; запазено е измерването.

\subsubsection{Клас 5 --- семантика на отстъпите и безопасните зони}

Класът, който най-ясно показва, че протичането не е между \emph{платформи}, а между
\emph{уговорки} --- две среди от едно и също семейство могат да се разминават помежду си.

\textbf{Разгледан пример --- двойно отчитане на безопасната зона.} Отстъпът на безопасната зона се
подава на съдържанието на страницата, но съдържанието вече е отместено от ръба и от собствения си
външен отстъп. На мобилната среда системната лента \emph{покрива} съдържанието --- двете величини
се мерят от един и същ ръб --- така че подаването на пълния отстъп върху вече отместено съдържание
го брои двукратно; поправката изважда собствения отстъп. Същата поправка обаче \emph{влоши}
настолната среда от същото семейство: там платформата на прозореца стои \emph{над} областта на
съдържанието, величините се сумират, и еталонът изобразява сбора им. Наложи се разклонение по
среда --- изваждане на едната, пълен отстъп на другата.

\textbf{Втори пример от същия клас.} Подравняването на съдържанието при списъчна контрола: еталонът
установява стойност, при която съдържанието минава под системната лента. Пренасянето ѝ поправи осем
интерфейса на мобилната среда и същевременно влоши настолната, където един и същ ред от изходния
текст на еталона изобразява различно.

Изводът е по-силен от двата случая поотделно: \emph{една и съща стойност в един и същ изходен текст
на еталона поражда различно изобразяване в две среди от едно семейство}. Абстракция, която
предполага, че семейството е еднородно, ще сгреши --- и ще сгреши безшумно, защото на едната среда
резултатът е правилен.

\subsubsection{Клас 6 --- поведение, изводимо само от реализацията на еталона}

Клас, при който разминаването не се дължи на уговорка, а на поведение, което не следва от описанието
и се открива само чрез прочитане на реализацията.

\textbf{Разгледан пример.} Заглавната и заключителната част на списъчна контрола се изобразяват или
не в зависимост от \emph{съчетанието} на вида разположение и дали елементите са групирани ---
вътрешното разпределение на еталона се разклонява по типа разположение и при групиране излиза
предсрочно, преди да достигне до тях. Системата беше приложила правилото на единия вид разположение
върху всички; поведението не е документирано никъде и не следва от нито един тест, изведено е само
от прочитане на разпределението по тип.

\subsubsection{Клас 7 --- разминавания, които не са дефекти}

Два подкласа, разграничени изрично:

\begin{enumerate}[topsep=2pt,itemsep=1pt,leftmargin=*]
\item \textbf{Особеност на конкретен платформен инструментариум}, която еталонът също търпи ---
      възпроизвеждането ѝ би било грешка.
\item \textbf{Артефакт на измерването} (§\ref{subsec:integrity}, §\ref{subsec:integrity_finding}).
\end{enumerate}

Класът съществува, защото и трите смущаващи фактора от §\ref{subsec:integrity} бяха първоначално
диагностицирани като дефекти на системата --- изброяването им предпазва следващия изследовател от
същата загуба.

\subsubsection{Изводи от систематизацията}

\textbf{Първо --- къде да се насочи вниманието.} Не към програмния интерфейс на контролите:
съответствието ,,свойство $\rightarrow$ native свойство“ се пренася почти без съпротива. Вниманието
принадлежи на конвенциите --- геометрични, стойностни, средови --- които никой инструментариум не
документира, защото за него самия те са очевидни.

\textbf{Второ --- разделение на средствата за откриване.} Класове 2, 3 и 6 са откриваеми без екран,
за сметка на преносимата логика, и принадлежат на модулното тестване. Класове 1, 4 и 5 изискват
сравнение на изображения, защото се проявяват само в състав или само на конкретна среда. Никое от
двете средства не покрива другото.

\textbf{Трето --- следствие за реда на работа.} Тъй като половината класове са невидими без
сравнение на изображения, измервателният апарат трябва да се изгражда \emph{успоредно} със
системата, а не след нея: изграден след това, той открива наведнъж натрупани разминавания с вече
преплетени причини; изграден успоредно, ги открива по едно, докато причината е още една.

\subsection{Цена на реализацията}
\label{subsec:cost}

Измерванията тук са възможни само защото предходните установяват \term{контролно условие}: двете
реализации изобразяват едно и също съдържание върху едни и същи native контроли; без него разлика в
размер или бързодействие не би била отличима от разлика в обхвата на реализираното.

Числата важат за платформи с потвърдено контролно условие, измерени с \code{measure\_size.py} върху
артефакти от 2026-08-16 и поправени за три дефекта в самия инструмент: несравними мерни състояния
(на диск срещу отрязано), пропуснато изключване на директории, създадени при изпълнение, и избор на
грешен главен файл при архивни формати. Единицата навсякъде по-долу е \textbf{MiB} ($2^{20}$~B);
стойности, публикувани по-рано с надпис ,,MB“, но вече изчислени в MiB, са преномерени само в
надписа.

\subsubsection{Размер на артефакта}

\begin{table}[H]
\caption{Размер на артефакта, отрязано срещу отрязано, release-компилация от двете страни (MiB, 2026-08-16)}
\label{tab:size}
\centering\small
\begin{tabular}{@{}lrrrl@{}}
\hline
\textbf{Платформа} & \textbf{еталон} & \textbf{система, по код} & \textbf{система, от опис.} & \textbf{съотношение (по код)} \\
\hline
Mac Catalyst        & 47,29  & 9,14  & 22,13 & 5,17$\times$ \\
iOS (симулатор)      & 49,64  & 9,21  & 22,52 & 5,39$\times$ \\
Windows\textsuperscript{1} & 149,18 & 9,17  & 10,38 & 16,3$\times$ \\
Android\textsuperscript{2}  & 34,33  & 30,56 & 36,86 & 1,12$\times$ (портът по-малък) \\
\hline
\end{tabular}
\end{table}

\footnotesize\textsuperscript{1} разгръщано състояние, без кеша на WebView2, записван при
изпълнение (виж (г) по-долу). \textsuperscript{2} една процесорна архитектура, разопаковано, изцяло
отрязано (виж (г) по-долу).\normalsize

\textbf{(а) Симетрия на обработката.} Чете се отрязаното състояние, не суровият размер на диска:
native артефактите се разпространяват без символи, затова диск и отрязано съвпадат, докато
управляваният носи 31{,}78 MiB символна информация в изпълнимия файл (Mac Catalyst) --- премахната
само от едната страна дава 8{,}61$\times$ вместо 5{,}17$\times$, завишение с 66\,\%.

\textbf{(б) Разпределение по компоненти (Mac Catalyst).} От 15,0 MiB метаданни и управлявани
библиотеки на еталона --- 8,53 MiB среда за изпълнение и базови библиотеки, 5,96 MiB код на
фреймуорка, 0,49 MiB код на приложението: средата за изпълнение води, но фреймуоркът е
\textbf{40\,\%} от разпределените байтове, не пренебрежим остатък. Преобладаващата компонента ---
31,5 MiB предварително компилиран образ --- остава \textbf{неразпределена}: маркерите в него са
структури от данни, не граници на код, а разбор на формата не е направен.

\textbf{(г) Три числа, публикувани по-рано, не издържаха повторна проверка} --- методологичен
резултат сам по себе си, не недостатък на порта, докладван честно вместо премълчан.

\emph{Android, ,,двойна архитектура“.} По-раншното твърдение, че портът не извлича собствената си
библиотека от APK-а, не се потвърждава: \code{extractNativeLibs} е \code{true} по подразбиране и
манифестът на порта не го отменя, така че платформата извлича библиотеката въпреки
STORED-опаковането. Реалната асиметрия е друга --- еталонният APK носи \emph{две} процесорни
архитектури (198 споделени библиотеки общо), портът --- \emph{една}; устройството зарежда само
една, затова сравнение с втората еталонна архитектура вътре сравнява две приложения с едно. При
еднакво условие съотношението за теглене се \emph{обръща}: 0{,}63$\times$, портът по-голям, а за
заемано, отрязано място е практически наравно, 1{,}12$\times$ (\tabref{tab:size}). Произходът е
конфигурационен от двете страни, не свойство на езика: еталонният проект не декларира архитектура и
изграждането взима подразбиращо и двете, докато портовият пресет избира изрично една.

\emph{Windows, кеш на WebView2.} Публикуваното по-рано съотношение 4,7$\times$ подценяваше
предимството на порта: 74{,}6\,\% от измерената директория на порта е кеш, записан от WebView2
\emph{при изпълнение}, не изграден артефакт. Изключен той (\tabref{tab:size}), съотношението е
16{,}3$\times$ --- най-голямото от четирите платформи, не най-малкото.

\emph{Windows, различен модел на разгръщане.} Еталонът носи \emph{цялата} self-contained среда
WindowsAppSDK със себе си (135{,}18 MiB, 90{,}5\,\% от товара му); портът разчита на инсталираната
на машината среда и носи само зареждащия компонент. Съотношението 16{,}3$\times$ измерва модел на
разгръщане, не ,,управляван срещу native“. Единственото еквивалентно сравнение --- собствен код
срещу собствен код --- дава 4,99 срещу 4,97 MiB, практически равни (с уговорката, че IL е по-плътен
от машинен код --- равенство по байтове, не по обем реализация).

\subsubsection{Защо декларативният вариант е по-обемен от кодовия}

Двата варианта се различават съществено по размер (\tabref{tab:size}). Естественото предположение
--- разликата идва от вградените описания --- е вярно само отчасти. Действителният обем на
вградения markup е \textbf{0,35 MiB} (не 0,78 MiB, изчислявано по-рано --- тази по-голяма стойност
всъщност е \emph{прирастът}, не обемът на самите описания).

\begin{table}[H]
\caption{Състав на два native артефакта на Mac Catalyst (release, без символна информация, MiB)}
\label{tab:sections}
\centering\small
\begin{tabular}{@{}lrrr@{}}
\hline
\textbf{Секция} & \textbf{по код} & \textbf{от описанието} & \textbf{разлика} \\
\hline
\code{\_\_text} (машинен код)      & 5,68 MiB & 9,65 MiB  & \textbf{+3,97 MiB} \\
\code{\_\_const} (константни данни) & 0,50 MiB & 1,29 MiB  & +0,79 MiB \\
\hline
\code{\_\_TEXT} общо               & 7,28 MiB & 14,34 MiB & +7,07 MiB \\
\hline
\end{tabular}
\end{table}

\footnotesize Секцията за низове \code{\_\_cstring} се различава с под 20~KiB между двата варианта
и не участва в анализа.\normalsize

Markup-ът присъства в артефакта \emph{двукратно} --- веднъж като вграден байтов масив, веднъж като
копие в константен низ, разбиран при първо извикване; удвоеният обем ($\approx$0,70 MiB) обяснява
86--90\,\% от прираста на \code{\_\_const} на трите платформи (macOS +0,79, iOS +0,81, Android
+0,78 MiB). Прирастът на \code{\_\_text} е около пет пъти по-голям --- разликата е машинен код, не
данни.

Причината следва от липсата на рефлексия (§\ref{subsec:xaml_choice}): кодовият вариант извиква само
конструкторите на контролите, които реално ползва, а свързващият редактор изхвърля останалото;
декларативният не знае предварително кой тип ще се появи и затова носи регистър на всички възможни
типове --- фабрика на обект и преобразувател ,,текст $\rightarrow$ стойност“ за всяко свойство ---
плюс разбора, обектния граф и разрешаването на свързванията. Всяка регистрация е код, недостъпен за
изхвърляне, защото се достига по данни; управляваната среда плаща същото, но веднъж, защото
механизмът ѝ за рефлексия вече е в средата за изпълнение.

На Android прирастът на кода е под една десета от Apple платформите (0,1 срещу $\approx$4,2 MiB):
портът там е споделена библиотека и всяка инстанциация на шаблон изнася манглирано име в таблицата
на символите --- прирастът се измества от код в изнесени символни имена. Причината за десеткратната
разлика по мащаб не е установена докрай (оптимизацията и броят на страниците са отхвърлени).

\subsubsection{Време до първи кадър (H2b-1)}

Единствената от трите хипотези за бързодействие, за която има измерване, по определението,
регистрирано предварително в §\ref{subsec:preregistration}: времето от стартиране до \emph{първия
заснет кадър, който не е начален екран} --- изрично \emph{не} до появата на прозорец. Всяка стойност
е медиана от 10 студени стартирания.

\input{generated/ttff_table}

Системата стартира по-бързо на двете измерими платформи: \textbf{1,42 пъти} на Android и
\textbf{1,73 пъти} на iOS. Двата входни пътя се държат еднакво --- декларативният не плаща за себе
си при стартиране, защото разборът е извършен при компилация (§\ref{subsec:xaml_impl}).

\textbf{Ограничение.} Разделителната способност на измерването е интервалът на извадката ---
около 0,24\,s. Разликите между двата входни пътя на системата (0,044\,s на Android, 0,062\,s на
iOS) са \emph{под} тази граница и не са разрешими --- докладват се като ,,еднакви“; разликата
спрямо еталона (0,43\,s и 0,84\,s) е няколкократно над границата и е разрешима. (По-ранен пробег на
същия инструмент даде за Android 3,1$\times$ вместо 1,42$\times$ при непроменен код --- разликата
идва от състоянието на емулатора; цитира се последният пълен пробег.)

Три платформи остават неизмерени по причина, записана в самата таблица: заснемането там минава през
виртуална машина и обиколката отнема от порядъка на самото измервано време.

\subsubsection{Потребена памет и останалите хипотези за времето}

Измерванията по \textbf{H2a} (потребена памет), \textbf{H2b-2} (разположение) и \textbf{H2b-3}
(установен режим) \emph{не са проведени} към момента на предаване. H2b-3 е формулирана \emph{срещу}
очакването на автора --- прогнозира липса на разлика при установен режим; потвърждаването на
H2b-1 не е довод в нейна полза, защото двете измерват различни неща.

\subsubsection{Съответствие с предварително регистрираните хипотези}

\begin{table}[H]
\caption{Съответствие между регистрираните хипотези и измереното}
\label{tab:hypotheses}
\centering\small
\begin{tabular}{@{}L{1.4cm}L{6.0cm}L{6.2cm}@{}}
\hline
\textbf{№} & \textbf{Прогноза (регистрирана предварително)} & \textbf{Изход} \\
\hline
H1 & управляваният артефакт е многократно по-обемен; разликата се дължи специфично на средата за
     изпълнение &
     \textbf{частично потвърдена} --- посоката се потвърждава навсякъде; величината варира силно:
     5,17--5,39$\times$ на Apple, 16,3$\times$ на Windows (собствен код 1,00$\times$), Android
     практически наравно по заемано място, по-голям за теглене (0,63$\times$) \\
H1' & разликата е голяма на Apple платформите и малка на Windows, където и двете колони изискват
      една и съща платформена среда &
     \textbf{частично потвърдена} --- вярна само за собствения код (1,00$\times$); еталонът на
     Windows е self-contained, не споделя средата, затова пълното съотношение е
     \emph{най-голямото} от всички (16,3$\times$) \\
H2a & по-малка следа на паметта & \textbf{непроверена} \\
H2b-1 & значително по-бързо студено стартиране & \textbf{потвърдена} --- 1,42$\times$ на Android,
        1,73$\times$ на iOS \\
H2b-2 & измеримо предимство при разположението & \textbf{непроверена} \\
H2b-3 & \emph{липса} на разлика при установен режим & \textbf{непроверена} \\
\hline
\end{tabular}
\end{table}

Три от шестте прогнози остават непроверени, таблицата го казва вместо да го премълчи. Ползата от
предварителната регистрация личи при H1 и H1': без нея формулировките биха били пренаписани
мълчешком след намирането на трите методологични дефекта в (г) по-горе; регистрираният текст остава
свидетелство какво точно е било предвидено, преди да е било измерено.

\subsection{Ограничения на изследването}
\label{subsec:limitations}

\begin{enumerate}[topsep=2pt,itemsep=1pt,leftmargin=*]
\item \textbf{Обхват на проверката за достоверност.} Механичната проверка (§\ref{subsec:integrity})
      улавя размножен кадър и неприложена тема, но не улавя кадър, заснет по време на преход, ако
      размерите му са правилни. Такъв кадър би понижил оценката, тоест грешката е в консервативната
      посока --- но остава неоткривана автоматично.
\item \textbf{Освобождаване на памет.} Изолацията на всяко измерване (§\ref{subsec:capture}) прави
      невъзможно наблюдението на освобождаване при напускане на интерфейс --- измерва се следа, не
      своевременност.
\item \textbf{Разделителна способност при измерване на стартиране} --- определението стъпва на
      честотата на заснемане и не претендира за милисекундна точност.
\item \textbf{Обвързаност с една версия на еталона и с едно измерване на размера.} Сравнението за
      степен на съвпадение е спрямо конкретна разпространена версия; числата за размера са от едно
      измерване (2026-08-16) и не се пренасят автоматично към следваща версия на нито едната страна.
\item \textbf{Подбор на интерфейсите.} Наборът е съставен за покритие на контролите и алгоритмите,
      не като извадка от типични приложения --- заключенията се отнасят за покритието, не за
      средното приложение.
\end{enumerate}
